% !TEX root = ../../../main.tex
% 来源：中学数学实验教材·第二册（几何）
% 章节：圆 / 圆与直线的位置关系
% 原始文件：4.tex，行 1932-1957
% 类型：tikz
% ---
\begin{tikzpicture}
\tkzDefPoint(-90:1.5){E}
\tkzDefPoint(0:1.5){F}
\tkzDefPoint(60:1.5){G}
\tkzDefPoint(150:1.5){H}
\tkzDefPoint(0,0){I}
\tkzDefTangent[at=E](I) \tkzGetPoint{a}
\tkzDefTangent[at=F](I) \tkzGetPoint{b}
\tkzDefTangent[at=G](I) \tkzGetPoint{c}
\tkzDefTangent[at=H](I) \tkzGetPoint{d}
\tkzInterLL(E,a)(F,b)\tkzGetPoint{B}
\tkzInterLL(F,b)(G,c)\tkzGetPoint{C}
\tkzInterLL(G,c)(H,d)\tkzGetPoint{D}
\tkzInterLL(E,a)(H,d)\tkzGetPoint{A}
\tkzDrawPolygon[thick](A,B,C,D)
\tkzDrawCircle[thick](I,E)
\tkzDrawSegments[dashed](I,E I,F I,G I,H I,D A,C B,I)
\tkzAutoLabelPoints[center=I](E,F,G,H,A,B,C,D)
\tkzInterLC(A,B)(B,C)  \tkzGetPoints{M}{M'}
\tkzInterLC(A,D)(D,C)  \tkzGetPoints{N}{N'}
\tkzDrawSegments[thick](M,N)
\tkzLabelPoints[below](M)
\tkzLabelPoints[left](N)
\tkzLabelPoints[above right](I)
\tkzDrawSegments[dashed](C,N C,M)
\end{tikzpicture}
